\section{Preprocessing Structure}
\subsection{Pipeline Definition}
Each preprocessing method is implemented as a deterministic function $f(\cdot)$ applied to waveform audio, producing a processed waveform:
\[
x' = f(x;\theta),
\]
where $\theta$ are fixed parameters. Pipelines are defined as ordered compositions (e.g., normalisation $\rightarrow$ enhancement $\rightarrow$ ASR), where normalisation is always applied as it is computationally cheap and necessary for consistent input to downstream models, but to avoid combinatorial explosion, only a limited set of combinations is evaluated.

\subsection{Experimental Plan}
The experimental study includes:
\begin{enumerate}[leftmargin=*]
    \item \textbf{Individual preprocessing methods:} baseline vs each preprocessing method in isolation.
    \item \textbf{Selected multi-step pipelines:} a small set of combinations chosen based on theoretical complementarity.
    \item \textbf{Speaker embedding experiment:} for target-speaker filtering, evaluate \textit{with} vs \textit{without} embeddings, holding all other factors fixed.
\end{enumerate}
If separation methods are constrained to $K=2$ due to model limitations, this is explicitly reported and separation results are presented as a focused subset.

\section{ASR Inference Controls}
To ensure fair comparisons, inference settings are held constant per model:
\begin{itemize}
    \item \textbf{Sampling rate:} audio is resampled to each model’s required input rate (typically 16\,kHz), with consistent resampling across conditions.
    \item \textbf{Decoding configuration:} beam size / temperature / language constraints are fixed and documented for each model.
    \item \textbf{Chunking/segmentation:} if models require chunked inference, chunk length and overlap are fixed and applied identically across conditions.
\end{itemize}
All model versions, checkpoints, and library commits are recorded for reproducibility.

\section{Evaluation Protocol}
\subsection{Reference Transcripts}
For synthetic mixtures, the reference includes the known transcript for each source speaker. For real datasets, references are taken from official annotations. Where multi-speaker references exist, meeting-style scoring is used (Section~\ref{sec:metrics_meeting}).

\subsection{Primary and Secondary Metrics}
\textbf{Primary outcome:} word error rate (WER) on the full clip reference(s).\\
\textbf{Secondary outcomes:} overlap-aware scores, meeting-style WER (where applicable), and error breakdowns. Metrics are defined formally in Section~\ref{sec:metrics}.

\subsection{Overlap-aware Evaluation}
\label{sec:metrics_overlap}
To directly answer RQ1 and RQ4, overlap-aware evaluation is performed by restricting scoring to words occurring within overlapped regions (using the overlap mask for synthetic mixtures, and overlap annotations/estimates for real audio). This yields \textbf{overlap-only WER}, which isolates performance under simultaneous speech from non-overlapped regions.

\subsection{Meeting-style Scoring for Multi-speaker References}
\label{sec:metrics_meeting}
When multi-speaker references are available and speaker attribution is not guaranteed in the hypothesis, meeting-style scoring such as concatenated permutation WER (cpWER) is used to reduce sensitivity to speaker ordering.
